\documentclass[tikz,border=8pt]{standalone}
\usepackage{ctex}
\usepackage{tikz}
\usetikzlibrary{arrows.meta,positioning,fit,backgrounds}

\begin{document}
\begin{tikzpicture}[
  font=\small,
  >=Latex,
  block/.style={draw=#1!70!black, fill=#1!8, rounded corners=2.5mm, very thick, minimum width=31mm, minimum height=11mm, align=center},
  flow/.style={->, very thick, draw=black!75},
  tip/.style={draw=black!35, fill=yellow!12, rounded corners=1.5mm, inner sep=3pt, align=center}
]

\node[block=blue] (raw) at (0,0) {原始 CAN 流\\CSV/TXT 报文序列};
\node[block=teal] (win) at (4.2,0) {滑窗构造\\\(W=64,\ S=32\)};
\node[block=teal] (feat) at (8.4,0) {特征编码\\8 字节 + \(\Delta t\)};
\node[block=purple] (cnn) at (12.6,0) {CAN-CNN(64x9)\\4 类 logits};
\node[block=green!70!black] (out) at (16.8,0) {部署输出\\class/confidence};

\draw[flow] (raw) -- (win);
\draw[flow] (win) -- (feat);
\draw[flow] (feat) -- (cnn);
\draw[flow] (cnn) -- (out);

\node[tip] at (8.4,-1.8) {关键点：\(\Delta t\) 仅用训练集统计归一化，\\部署端复用同一 \texttt{dt\_max\_ms} 保证分布一致性。};

\begin{scope}[on background layer]
  \node[draw=black!25, rounded corners=3mm, line width=0.8pt, fit=(raw)(win)(feat)(cnn)(out), inner sep=7mm,
  label={[yshift=2mm]\large\bfseries CAN-CNN(64x9) 方法总览（重绘版）}] {};
\end{scope}
\end{tikzpicture}
\end{document}
